% from ega4-9-fr.tex
\oldpage[IV-3]{56}
propriétés envisagées dans (8.10.5) sont toutes stables par changement de base.
La condition (iii) de (9.1.1) est donc satisfaite en vertu de (8.8.2, (i))
(resp. (8.10.5)), puisque $\operatorname{Spec}(k)$ est quasi-compact et
quasi-séparé. Reste à vérifier la condition (ii) de (9.1.1) qui découle encore
de (8.8.2, (i)) et de (8.10.5), en considérant $K$ comme limite inductive de
ses sous-extensions de type fini. On conclut donc par (9.1.1).

En particulier, s'il existe une extension $K$ de $k$ et un $K$-isomorphisme
$X_{(K)}\xrightarrow{\sim}Y_{(K)}$, on dit que $X$ et $Y$ sont
\emph{géométriquement isomorphes}.

Le corollaire suivant généralise (II, 6.6.5)~:

\begin{corollary}[9.1.5]
\label{IV.9.1.5-fr}
Soient $k$ un corps, $X$ un $k$-préschéma. S'il existe une extension $K$ de
$k$ telle que $X_{(K)}$ soit projectif (resp. quasi-projectif) sur $K$, alors
$X$ est projectif (resp. quasi-projectif) sur $k$.
\end{corollary}

Le morphisme $\operatorname{Spec}(K)\to\operatorname{Spec}(k)$ étant
fidèlement plat et quasi-compact, il résulte déjà de (2.7.1, (v)) que $X$ est
de type fini sur $k$. L'hypothèse signifie qu'il existe une immersion fermée
(resp. une immersion) $X_{(K)}\to\mathbf P_K^r=\mathbf P_k^r\otimes_k K$
(II, 5.5.4, (ii) et 5.5.3)~; appliquant (9.1.4) pour la propriété (iv)
(resp. (ii)) de (8.10.5), on en déduit qu'il y a une extension finie $k'$ de
$k$ et une immersion fermée (resp. une immersion)
$X_{(k')}\to\mathbf P_{k'}^r$, autrement dit $X_{(k')}$ est projectif
(resp. quasi-projectif) sur $k'$. On conclut alors par (II, 6.6.5).

\begin{proposition}[9.1.6]
\label{IV.9.1.6-fr}
Soient $k$ un corps, $\Omega$ une extension algébriquement close de $k$,
$X$ un préschéma de type fini sur $k$, $\mathcal F$, $\mathcal G$ deux
$\mathcal O_X$-Modules cohérents. Supposons qu'il existe un isomorphisme
$\mathcal F\otimes_k\Omega\xrightarrow{\sim}\mathcal G\otimes_k\Omega$.
Alors il existe une extension finie $k'$ de $k$ et un isomorphisme
$\mathcal F\otimes_k k'\xrightarrow{\sim}\mathcal G\otimes_k k'$.
\end{proposition}

Le raisonnement est le même que dans (9.1.4), en appliquant (8.5.2, (i))
(on utilise ici, dans la démonstration de la propriété (iii) de (9.1.1), le
fait que les morphismes $D(g)\to D(f)$ (avec les notations de (9.1.3)) sont
des immersions ouvertes, et \emph{a fortiori} des morphismes plats).

\subsection{Propriétés constructibles et ind-constructibles}
\label{subsection:IV.9.2-fr}

\begin{definition}[9.2.1]
\label{IV.9.2.1-fr}
Soit $\mathbf P(X,\mathcal F,k)$ une relation. On dit que $\mathbf P$ est une
propriété \emph{constructible} (resp. \emph{ind-constructible}) de préschémas
algébriques si les deux conditions suivantes sont vérifiées~:
\begin{enumerate}[label=(\roman*)]
  \item Si $k$ est un corps, $X$ un préschéma algébrique sur $k$,
  $\mathcal F$ un $\mathcal O_X$-Module cohérent, $k'$ une extension de $k$,
  alors, pour que $\mathbf P(X,\mathcal F,k)$ soit vraie, il faut et il suffit
  que $\mathbf P(X_{(k')},\mathcal F\otimes_k k',k')$ soit vraie.
  \item Soient $S$ un préschéma intègre noethérien, de point générique $\eta$,
  $u:X\to S$ un morphisme de type fini, $\mathcal F$ un $\mathcal O_X$-Module
  cohérent. Pour tout $s\in S$, posons
  $X_s=u^{-1}(s)=X\times_S\operatorname{Spec}(k(s))$,
  $\mathcal F_s=\mathcal F\otimes_{\mathcal O_S}k(s)$. Soit $E$ l'ensemble
  des $s\in S$ tels que $\mathbf P(X_s,\mathcal F_s,k(s))$ soit vraie.
  Alors l'un des ensembles $E$, $S-E$ (resp. l'ensemble $E$) contient un
  ouvert non vide (et par suite est un voisinage de $\eta$) (resp. contient
  un ouvert non vide s'il contient $\eta$).
\end{enumerate}
\end{definition}

